\section{Industrial context}
\label{sec:context}

This section is context, not evidence. No claim in \S\ref{sec:evaluation} rests on anything reported here, and the material is included because it explains why the question was asked at all and why the artifact is small.

The work began inside FlowForge, a larger typed pipeline framework. Two things learned there shaped this paper. The first is that a typed example can drift away from the execution path it appears to describe: a framework can present a convincing compile-time story while the frame actually written reaches the sink by a different route, so the demonstration and the checked boundary are not the same place. That is why the artifact here ties the compile-time witness to the sink-construction call rather than to a separate demonstration, and why the runtime pin sits on the write path. The second is that a framework narrative accumulates connectors, evolution rules, observability and positioning long before its enforcement path is closed, which makes it a poor subject for a claim about enforcement.

Both lessons are about overclaiming, and the same failure mode showed up inside this paper's own measurements. The harness was built to characterise Spark and the first thing it characterised was the artifact: the runtime comparator read two carriers strictly and the third not at all (\S\ref{sec:weakness}). A framework-scale evaluation would not have found that, because the defect is invisible unless every carrier is varied independently against every predicate. The decision to extract a small artifact and measure it directly is therefore not modesty about scope; it is what made the result observable.
